% i18n_qa_ro_paper.tex — Romanian engineering paper on signal processing
% Intentional errors: cedilla (ş/ţ) instead of comma-below (ș/ț),
%   wrong Romanian quotes, encoding issues
\documentclass[12pt,a4paper]{article}

\usepackage[romanian]{babel}
\usepackage[utf8]{inputenc}
\usepackage[T1]{fontenc}
\usepackage{amsmath}
\usepackage{amssymb}
\usepackage{graphicx}
\usepackage{hyperref}
\usepackage[margin=2.5cm]{geometry}
\usepackage{booktabs}
\usepackage{natbib}
\usepackage{enumitem}
\usepackage{algorithm}
\usepackage{algpseudocode}
\usepackage{microtype}
\usepackage{xcolor}
\usepackage{float}
\usepackage{subcaption}
\usepackage{siunitx}

% INTENTIONAL ERROR: cedilla ş instead of comma-below ș in title
\title{Metode avansate de prelucrare a semnalelor digitale\\
pentru sisteme de comunicaţii 5G}
\author{%
  Alexandru Ionescu\thanks{Departamentul de Electronică şi Telecomunicaţii,
  Universitatea Politehnica din Bucureşti.}
  \and
  Maria Popescu\thanks{Institutul Naţional de Cercetare-Dezvoltare
  pentru Microtehnologie, Bucureşti.}
}
\date{6 aprilie 2026}

\begin{document}

\maketitle

\begin{abstract}
% INTENTIONAL ERROR: cedilla ş/ţ throughout instead of comma-below ș/ț
Acest articol prezintă o metodă nouă de prelucrare a semnalelor digitale
bazată pe transformata wavelet discretă şi algoritmi de comprimare adaptivă.
Propunem un cadru teoretic unificat pentru analiza spectrală a semnalelor
în sisteme de comunicaţii de generaţia a cincea (5G). Metoda noastră
permite reducerea complexităţii computaţionale cu 35\% faţă de abordările
convenţionale bazate pe transformata Fourier rapidă. Rezultatele
experimentale demonstrează că algoritmul propus atinge un raport semnal-zgomot
superior cu \SI{4.2}{\decibel} în condiţii de canal cu estompare Rayleigh.
Validarea a fost realizată atât prin simulări Monte Carlo, cât şi prin
măsurători pe o platformă hardware FPGA Xilinx Zynq UltraScale+.
\end{abstract}

\tableofcontents

%% ---------------------------------------------------------------
\section{Introducere}
\label{sec:intro}

Sistemele de comunicaţii de generaţia a cincea (5G) necesită tehnici
avansate de prelucrare a semnalelor pentru a gestiona lăţimile de bandă
extinse şi densităţile spectrale ridicate. Standardul 3GPP Release~17
% INTENTIONAL ERROR: wrong quotes — English "..." instead of Romanian „..."
specifică "multiple-input multiple-output" (MIMO) masiv cu până la
256~de antene la staţia de bază~\citep{dahlman2020}.

Transformata Fourier rapidă (FFT) a fost instrumentul fundamental în
prelucrarea semnalelor digitale timp de decenii. Cu toate acestea,
pentru semnale nestaţionare întâlnite în canalele 5G cu mobilitate
ridicată, analiza timp-frecvenţă oferită de transformata wavelet
discretă (DWT) este superioară~\citep{mallat2009}.

% INTENTIONAL ERROR: cedilla ţ instead of comma-below ț
Contribuţiile principale ale acestui articol sunt:
\begin{enumerate}[label=(\roman*)]
  \item Un algoritm de descompunere wavelet adaptivă cu complexitate
        $\mathcal{O}(N \log N)$ în loc de $\mathcal{O}(N^2)$;
  \item O metodă de estimare a canalului bazată pe proiecţii pe
        subraze wavelet;
  \item Validare experimentală pe platforma FPGA cu latentă sub
        \SI{1}{\milli\second}.
\end{enumerate}

Restul articolului este organizat astfel: Secţiunea~\ref{sec:fundamente}
prezintă fundamentele teoretice, Secţiunea~\ref{sec:metoda} descrie
metoda propusă, Secţiunea~\ref{sec:rezultate} conţine rezultatele
experimentale, iar Secţiunea~\ref{sec:concluzii} prezintă concluziile.

%% ---------------------------------------------------------------
\section{Fundamente teoretice}
\label{sec:fundamente}

\subsection{Transformata Fourier discretă}

Pentru un semnal discret $x[n]$ de lungime $N$, transformata Fourier
discretă (DFT) este definită ca:
\begin{equation}
\label{eq:dft}
X[k] = \sum_{n=0}^{N-1} x[n] \, e^{-j 2\pi kn/N}, \quad k = 0, 1, \ldots, N-1.
\end{equation}

Algoritmul FFT calculează DFT cu complexitate $\mathcal{O}(N \log N)$
prin descompunere recursivă\footnote{Algoritmul Cooley--Tukey este cea
mai utilizată variantă, publicat în 1965.}. Densitatea spectrală de
putere se obţine din:
\begin{equation}
\label{eq:psd}
S_{xx}[k] = \frac{1}{N} |X[k]|^2.
\end{equation}

\subsection{Transformata wavelet discretă}

Transformata wavelet discretă descompune semnalul pe niveluri de
rezoluţie folosind filtre trece-jos şi trece-sus:
\begin{align}
\label{eq:dwt_low}
a_j[n] &= \sum_k h[k-2n] \, a_{j-1}[k], \\
\label{eq:dwt_high}
d_j[n] &= \sum_k g[k-2n] \, a_{j-1}[k],
\end{align}
unde $h[n]$ şi $g[n]$ sunt filtrele de analiză wavelet, $a_j[n]$ sunt
coeficienţii de aproximare, iar $d_j[n]$ sunt coeficienţii de detaliu
la nivelul $j$.

% INTENTIONAL ERROR: wrong Romanian quotes
Condiţia de "reconstucţie perfectă" cere ca filtrele să satisfacă:
\begin{equation}
\label{eq:pr}
H(z)H^*(z^{-1}) + H(-z)H^*(-z^{-1}) = 2,
\end{equation}
unde $H(z)$ este funcţia de transfer a filtrului trece-jos.

\subsection{Modelul canalului 5G}

Modelul de canal pentru sisteme 5G cu frecvenţe milimetrice include
propagare cu raze multiple:
\begin{equation}
\label{eq:canal}
h(\tau, t) = \sum_{l=0}^{L-1} \alpha_l(t) \, \delta(\tau - \tau_l(t)),
\end{equation}
unde $\alpha_l(t)$ este câştigul complex al căii $l$, iar $\tau_l(t)$
este întârzierea corespunzătoare. Raportul semnal-zgomot (SNR) este:
\[
\mathrm{SNR} = \frac{E_s}{N_0} = \frac{P_s}{B \cdot N_0},
\]
unde $P_s$ este puterea semnalului, $B$ este lăţimea de bandă, iar
$N_0$ este densitatea spectrală de putere a zgomotului.

%% ---------------------------------------------------------------
\section{Metoda propusă}
\label{sec:metoda}

\subsection{Descompunere wavelet adaptivă}

Propunem un algoritm de descompunere wavelet care adaptează numărul
de niveluri în funcţie de caracteristicile spectrale ale semnalului
recepţionat. Fie $x[n]$ semnalul recepţionat şi $J_{\max}$ numărul
maxim de niveluri de descompunere.

\begin{definition}
Criteriul de oprire adaptivă este definit ca:
\begin{equation}
\label{eq:criteriu}
\mathcal{C}_j = \frac{\|d_j\|_2^2}{\|a_{j-1}\|_2^2} < \epsilon,
\end{equation}
unde $\epsilon$ este un prag predefinit dependent de SNR-ul estimat.
\end{definition}

Algoritmul procedează astfel:

\begin{algorithm}[ht]
\caption{Descompunere wavelet adaptivă}
\label{alg:dwt_adaptiv}
\begin{algorithmic}[1]
\Require Semnal $x[n]$, prag $\epsilon$, nivel maxim $J_{\max}$
\Ensure Coeficienţi $\{a_J, d_J, d_{J-1}, \ldots, d_1\}$
\State $a_0 \gets x$
\For{$j = 1$ to $J_{\max}$}
  \State Calculează $a_j, d_j$ conform (\ref{eq:dwt_low}), (\ref{eq:dwt_high})
  \State Calculează $\mathcal{C}_j$ conform (\ref{eq:criteriu})
  \If{$\mathcal{C}_j < \epsilon$}
    \State $J \gets j$
    \State \textbf{break}
  \EndIf
\EndFor
\State \Return $\{a_J, d_J, d_{J-1}, \ldots, d_1\}$
\end{algorithmic}
\end{algorithm}

\subsection{Estimarea canalului prin proiecţii wavelet}

% INTENTIONAL ERROR: cedilla ş instead of comma-below ș
Estimarea canalului se realizează prin proiecţia semnalului pilot
pe un spaţiu wavelet ales adaptiv. Considerăm semnalul pilot
$p[n]$ transmis pe subpurtătoarele OFDM şi semnalul recepţionat
$y[n] = h[n] * p[n] + w[n]$, unde $w[n]$ este zgomot gaussian alb.

Estimatorul propus este:
\begin{equation}
\label{eq:estimator}
\hat{H}[k] = \sum_{j=1}^{J} \sum_{m} \hat{d}_{j,m} \, \psi_{j,m}[k] + \sum_{m} \hat{a}_{J,m} \, \phi_{J,m}[k],
\end{equation}
unde $\psi_{j,m}$ şi $\phi_{J,m}$ sunt funcţiile wavelet şi de scalare,
iar coeficienţii sunt estimaţi prin:
\[
\hat{d}_{j,m} = \begin{cases}
\tilde{d}_{j,m} & \text{dacă } |\tilde{d}_{j,m}| > \lambda_j, \\
0 & \text{altfel},
\end{cases}
\]
cu pragul de tăiere $\lambda_j = \sigma_w \sqrt{2 \ln N / 2^j}$.

\subsection{Complexitatea computaţională}

\begin{table}[ht]
\centering
% INTENTIONAL ERROR: wrong quotes again
\caption{Comparaţie a complexităţii computaţionale între metodele "clasice" şi metoda propusă}
\label{tab:complexitate}
\begin{tabular}{lccc}
\toprule
Metodă & Adunări & Înmulţiri & Total \\
\midrule
FFT standard & $N \log_2 N$ & $\frac{N}{2} \log_2 N$ & $\mathcal{O}(N \log N)$ \\
DWT complet ($J$ niveluri) & $2NL$ & $2NL$ & $\mathcal{O}(NL)$ \\
DWT adaptiv (propus) & $2N'L$ & $2N'L$ & $\mathcal{O}(N'L)$ \\
MMSE clasic & $N^3$ & $N^3$ & $\mathcal{O}(N^3)$ \\
LS clasic & $N^2$ & $N^2$ & $\mathcal{O}(N^2)$ \\
Wavelet-LS (propus) & $N \log N$ & $N \log N$ & $\mathcal{O}(N \log N)$ \\
\bottomrule
\end{tabular}
\end{table}

În Tabelul~\ref{tab:complexitate}, $N' \leq N$ este dimensiunea efectivă
determinată de criteriul de oprire~\eqref{eq:criteriu}, iar $L$ este
lungimea filtrului wavelet. Pentru wavelet-ul Daubechies db4, $L = 8$.

%% ---------------------------------------------------------------
\section{Rezultate experimentale}
\label{sec:rezultate}

\subsection{Configuraţia simulărilor}

Simulările au fost realizate în MATLAB R2025a pe un sistem cu
procesor Intel Core i9-14900K. Parametrii sistemului OFDM sunt
prezentaţi în continuare\footnote{Parametrii corespund standardului
3GPP TS~38.211 pentru banda de frecvenţe FR2 (24--52~GHz).}:
\begin{itemize}
  \item Număr de subpurtătoare: $N = 2048$;
  \item Spaţiere subpurtătoare: $\Delta f = \SI{120}{\kilo\hertz}$;
  \item Prefix ciclic: $N_{\mathrm{CP}} = 144$;
  \item Modulaţie: 64-QAM;
  \item Model canal: CDL-C (clustered delay line)~\citep{3gpp2020}.
\end{itemize}

\subsection{Performanţa estimării canalului}

Figura~\ref{fig:mse} prezintă eroarea medie pătratică (MSE) a estimării
canalului în funcţie de SNR pentru diferite metode.

\begin{figure}[ht]
\centering
\fbox{\parbox{0.75\textwidth}{%
\centering
\textbf{MSE vs. SNR pentru estimare canal CDL-C}\\[8pt]
\begin{tabular}{cc}
SNR (dB) & MSE (propus / LS / MMSE) \\
\hline
0 & $10^{-1.2}$ / $10^{-0.8}$ / $10^{-1.3}$ \\
10 & $10^{-2.5}$ / $10^{-1.8}$ / $10^{-2.4}$ \\
20 & $10^{-3.8}$ / $10^{-2.8}$ / $10^{-3.5}$ \\
30 & $10^{-4.2}$ / $10^{-3.0}$ / $10^{-4.0}$ \\
\end{tabular}\\[6pt]
\emph{Metoda propusă (linie continuă roşie) depăşeşte}\\
\emph{LS şi este comparabilă cu MMSE la SNR ridicat.}
}}
\caption{Eroarea medie pătratică a estimării canalului în funcţie de raportul
semnal-zgomot. Canalul CDL-C cu viteză \SI{120}{\kilo\meter\per\hour}.}
\label{fig:mse}
\end{figure}

Rezultatele arată că metoda propusă atinge o performanţă superioară
estimatorului LS cu \SI{10}{\decibel} la SNR = \SI{20}{\decibel} şi
este comparabilă cu estimatorul MMSE, dar cu complexitate
semnificativ mai redusă.

\subsection{Rata de eroare pe biţi}

Rata de eroare pe biţi (BER) pentru modulaţie 64-QAM este prezentată
în Figura~\ref{fig:ber}.

\begin{figure}[ht]
\centering
\fbox{\parbox{0.75\textwidth}{%
\centering
\textbf{BER vs. $E_b/N_0$ pentru 64-QAM, canal CDL-C}\\[8pt]
\begin{tabular}{cc}
$E_b/N_0$ (dB) & BER (propus / LS / perfect) \\
\hline
5 & $3.2 \times 10^{-2}$ / $5.1 \times 10^{-2}$ / $2.8 \times 10^{-2}$ \\
10 & $4.7 \times 10^{-3}$ / $1.2 \times 10^{-2}$ / $3.9 \times 10^{-3}$ \\
15 & $2.1 \times 10^{-4}$ / $8.3 \times 10^{-3}$ / $1.5 \times 10^{-4}$ \\
20 & $6.8 \times 10^{-6}$ / $2.1 \times 10^{-3}$ / $4.2 \times 10^{-6}$ \\
\end{tabular}\\[6pt]
\emph{Diferenţa faţă de estimare perfectă: sub \SI{1}{\decibel}.}
}}
\caption{Rata de eroare pe biţi în funcţie de $E_b/N_0$ pentru modulaţie 64-QAM.
Comparaţie între metoda propusă, estimare LS şi cunoaştere perfectă a canalului.}
\label{fig:ber}
\end{figure}

\subsection{Implementare FPGA}

% INTENTIONAL ERROR: cedilla ţ instead of comma-below ț
Implementarea pe FPGA Xilinx Zynq UltraScale+ ZCU111 a fost realizată
în Vivado HLS 2025.1. Resursele utilizate sunt prezentate în
Tabelul~\ref{tab:fpga}.

\begin{table}[ht]
\centering
\caption{Utilizarea resurselor FPGA pentru algoritmul propus}
\label{tab:fpga}
\begin{tabular}{lrrc}
\toprule
Resursă & Utilizat & Disponibil & Procent \\
\midrule
LUT & 45\,230 & 425\,280 & 10.6\% \\
FF & 62\,480 & 850\,560 & 7.3\% \\
BRAM & 128 & 1\,080 & 11.9\% \\
DSP & 256 & 4\,272 & 6.0\% \\
\bottomrule
\end{tabular}
\end{table}

Latenţa măsurată pentru o subramă OFDM completă ($N = 2048$)
este de \SI{0.82}{\milli\second}, sub pragul de \SI{1}{\milli\second}
specificat în cerinţele 5G NR\footnote{%
Specificaţia 3GPP TS~38.214 v17.3.0, Secţiunea~5.1.}.

%% ---------------------------------------------------------------
\section{Discuţii}
\label{sec:discutii}

\subsection{Alegerea wavelet-ului}

Am comparat performanţele pentru mai multe familii wavelet:
Daubechies (db2--db10), Symlet (sym4--sym8) şi Coiflet (coif2--coif5).
Cel mai bun compromis între performanţă şi complexitate a fost obţinut
cu wavelet-ul Daubechies db4, care are 8~coeficienţi de filtru şi
4~momente nule. Rezultatele sunt consistente cu analiza teoretică
din~\citet{mallat2009}.

\subsection{Comparaţie cu metode de învăţare profundă}

Recent, metodele bazate pe reţele neuronale au fost propuse pentru
estimarea canalului~\citep{ye2018}. Deşi aceste metode pot atinge
performanţe superioare în condiţii staţionare, ele necesită
antrenament extensiv şi nu generalizează bine la scenarii noi.
Metoda noastră bazată pe wavelet este complet analitică şi nu necesită
date de antrenament.

\subsection{Limitări}

Principala limitare a metodei propuse constă în presupunerea că
canalul are o reprezentare rară în domeniul wavelet. Pentru canale
cu difuzie bogată (de exemplu, medii interioare dense), această
presupunere poate fi violată, iar performanţa se degradează.

%% ---------------------------------------------------------------
\section{Concluzii}
\label{sec:concluzii}

Am prezentat o metodă nouă de estimare a canalului pentru sisteme
5G bazată pe transformata wavelet discretă adaptivă. Metoda reduce
complexitatea computaţională cu 35\% faţă de abordările
convenţionale şi atinge un câştig de \SI{4.2}{\decibel} în SNR faţă
de estimarea LS. Implementarea FPGA demonstrează fezabilitatea
metodei pentru aplicaţii în timp real, cu latenţă sub
\SI{1}{\milli\second}. Lucrările viitoare vor investiga extensia
metodei la sisteme MIMO masiv şi la scenarii cu mobilitate
ultra-ridicată.

%% ---------------------------------------------------------------
\begin{thebibliography}{99}

\bibitem[3GPP(2020)]{3gpp2020}
3GPP (2020).
TR~38.901 v16.1.0: Study on channel model for frequencies from 0.5 to 100~GHz.
\emph{3rd Generation Partnership Project}.

\bibitem[Dahlman et al.(2020)]{dahlman2020}
Dahlman, E., Parkvall, S. \c{s}i Sk\"{o}ld, J. (2020).
\emph{5G NR: The Next Generation Wireless Access Technology}, ed.~a 2-a.
Academic Press.

\bibitem[Mallat(2009)]{mallat2009}
Mallat, S. (2009).
\emph{A Wavelet Tour of Signal Processing: The Sparse Way}, ed.~a 3-a.
Academic Press.

\bibitem[Proakis \c{s}i Manolakis(2014)]{proakis2014}
Proakis, J.\,G. \c{s}i Manolakis, D.\,G. (2014).
\emph{Digital Signal Processing: Principles, Algorithms, and Applications},
ed.~a 4-a. Pearson.

\bibitem[Ye et al.(2018)]{ye2018}
Ye, H., Li, G.\,Y. \c{s}i Juang, B.-H. (2018).
Power of deep learning for channel estimation and signal detection in OFDM systems.
\emph{IEEE Wireless Communications Letters}, 7(1), 114--117.

\bibitem[Tse \c{s}i Viswanath(2005)]{tse2005}
Tse, D. \c{s}i Viswanath, P. (2005).
\emph{Fundamentals of Wireless Communications}.
Cambridge University Press.

\end{thebibliography}

\end{document}
